\documentclass[11pt]{article}
\usepackage{amsmath}
\usepackage{amsfonts}
\usepackage{amsthm}
\usepackage[utf8]{inputenc}
\usepackage[margin=0.75in]{geometry}

\title{CSC110 Fall 2026 Assignment 3: Number Theory and Cryptography}
\author{Idzudin Saffuan}
\date{\today}
\begin{document}

\maketitle


\section*{Part 1: Proofs}

\begin{enumerate}
\item[1.] Statement to prove:
$\forall a, b, k, m, n \in \mathbb{Z},~
\big(n > 0 \land a \equiv b \pmod n \big) \Rightarrow
a + kn \equiv b + mn \pmod n$

\begin{proof}
Let $a, b, k, m, n \in \mathbb{Z}$ be arbitrary. Suppose $n > 0 \wedge a \equiv b \ (mod\ n)$.

Then $n |(a-b) \implies a - b = nk_1$ for some $k \in \mathbb{Z}$. Consider $a +kn$ and $b + mn$:
\begin{align*}
a+kn - (b +mn) \\
= a-b + kn - mn \\
= a-b +(k-m)n
\end{align*}
Notice that $a-b = nk_1$ and $k, m \in \mathbb{Z}$. Then:
\begin{align*}
n(k_1+k-m) \implies n|\big((a+kn)-(b+mn)\big)
\end{align*}
Hence $a + kn \equiv b + mn \ (mod\ n)$ as desired.
\end{proof}

\newpage

\item[2.] Statement to prove:
$\forall a, b, n \in \mathbb{Z},~
\big(n > 0 \land \gcd(a, n) = 1 \land b \equiv a \pmod n \big) \Rightarrow
\gcd(b, n) = 1$

\begin{proof}
Let $a, b, n \in \mathbb{Z}$ be arbitrary. Suppose $n>0 \wedge gcd(a,n) =1$ and $ b \equiv a \ (mod \ n)$.

By the \textit{GCD Linear Combination Theorem},
\begin{align*}
pa + qn = 1
\end{align*}
for some $p, q \in \mathbb{Z}$. Recall that $ b \equiv a \ (mod \ n)$ then by definition:
\begin{align*}
n|(b-a) \implies b-a = nk
\end{align*}
for some $k \in \mathbb{Z}$. We know $ a= b-nk$;
\begin{align*}
p(b-nk) + qn = 1 \\
\implies pb - pnk + qn = 1 \\
\implies pb + (q-pk)n = 1
\end{align*}
We take $x=p$, $y = q-pk$ such that $x,y \in \mathbb{Z}$. Hence, it is shown that there exists a linear combination $xb +yn =1$. Recall that $n > 0$ and $n,b \in \mathbb{Z}$ such that $gcd(b,n) \neq  0$. Since 1 is the smallest positive integer that could be written as a linear combination of $b$ and $n$, It follows that $gcd(b,n) = 1$
\end{proof}

\newpage

\item[3.] Statement to prove:
Let $p$ and $q$ be distinct primes, let $n = pq$, and let $\varphi(n) = (p - 1)(q - 1)$.
Let $e, d \in \mathbb{Z}^+$ and assume that $ed \equiv 1 \pmod{\varphi(n)}$.
Prove that for all $x \in \mathbb{Z}$, if $\gcd(x, n) = 1$, then
\[
(x^e)^d \equiv x \pmod n.
\]

\begin{proof}
Suppose $gcd(x,n) =1$. Then by \textit{Euler's Theorem}:
\begin{align*}
x^{\phi(n)} \equiv 1 \ (mod \ n)
\end{align*}
Since $ed \equiv 1 \ \big(mod \ \phi(n)\big)$, by definition:
\begin{align*}
\phi(n)|(ed-1) \\
\implies k\phi(n) = ed - 1 \\
\implies k\phi(n) +1 = ed
\end{align*}
Consider:
\begin{align*}
x^{\phi(n)} \equiv 1 \ (mod \ n) \\
=(x^{\phi(n)})^k \equiv 1^k \ (mod \ n) \\
=x^{k\phi(n)} \equiv 1 \ (mod \ n) \\
=x \cdot x^{k\phi(n)} \equiv x \cdot 1 \ (mod \ n) \\
=x^{k\phi(n)+1} \equiv x\ (mod \ n) \\
\end{align*}
Recall that $k\phi(n) +1 = ed$:
\begin{align*}
x^{ed} \equiv x\ (mod \ n) \\
= (x^{e})^d \equiv x\ (mod \ n) \\
\end{align*}
as desired.
\end{proof}
\end{enumerate}


\newpage

\section*{Part 2: Block-Based RSA}

Complete this part in the provided \texttt{a3\_part2.py} starter file.
Do \textbf{not} include your solutions in this file.

\section*{Part 3: Digital Signatures}

\subsection*{Part (a): Introduction}

Complete this part in the provided \texttt{a3\_part3.py} starter file.
Do \textbf{not} include your solutions in this file.

\subsection*{Part (b): Generalizing the message digests}

Complete most of this part in the provided \texttt{a3\_part3.py} starter file.
Do \textbf{not} include your solutions in this file, \emph{except} for the following two questions:

\begin{enumerate}

\item[3b.]

\begin{verbatim}


def find_collision_len_times_sum(message: str) -> str:
new_message_so_far = ''
    first = chr(ord(message[0]) - 1)
    new_message_so_far = new_message_so_far + first
    for i in range(1, len(message) - 1):
        new_message_so_far = new_message_so_far + message[i]
    last = chr(ord(message[len(message) - 1]) + 1)
    new_message_so_far = new_message_so_far + last
    return new_message_so_far

\end{verbatim}


The function calculates the sum of the \texttt{ord} values of the strings within \texttt{message}. Substracting $1$ from the \texttt{ord} value of the first character in \texttt{message} and adding $1$ in last will give the same sum as the original message. Hence the signature generated using the new message will be the same as the original message.

\newpage

\item[4b.]

\begin{verbatim}


def find_collision_ascii_to_int(public_key: tuple[int, int], message: str) -> str:
n = public_key[0]
    int_128 = ascii_to_int(message) + n
    num = a3_part2.int_to_base128(int_128)
    str_so_far = ''
    for i in num:
        str_so_far = str_so_far + chr(i)
    return str_so_far
\end{verbatim}

The message digest in this case is calculated by finding the integer value from the given \texttt{message}'s \texttt{ord} values in base-128. Recall that the signature is calculated by:
\begin{align*}
s = a^d  \ \% \ n
\end{align*}
where $a$ is the message digest modulo $n$, and a is calculated from the return value of \texttt{ascii\_to\_int} modulo $n$. From \textit{Part 1, Question 1}, We know that:
\begin{align*}
a + kn \equiv b + mn \ (mod\ n)
\end{align*}
The integer value added by $n$ is equivalent to the original integer value modulo $n$, from which the signature is generated from.
\end{enumerate}
\end{document}
